\begin{workedexample}[Worked Example: SAR scaling with field]
Specific absorption rate is the RF power deposited per unit mass. Because
heating goes as the square of the RF field and frequency,
\[ \text{SAR} \propto f^2 B_1^2 \propto B_0^2\,\alpha^2. \]
Moving from $1.5\ \text{T}$ to $3\ \text{T}$ at the same flip angle raises SAR about
fourfold ($f$ doubles). The IEC caps whole-body SAR (normal mode $2\ \text{W/kg}$),
so high-field sequences must lengthen TR, lower flip angles, or use low-SAR
refocusing to stay within limits.
\end{workedexample}

\begin{keytakeaways}
\begin{itemize}[leftmargin=1.4em,itemsep=2pt]
\item $\text{SAR} \propto f^2 B_1^2$, so it scales as $B_0^2$ for a given flip angle.
\item IEC normal-mode whole-body limit is $2\ \text{W/kg}$; first-level mode allows more under supervision.
\item High field forces SAR management: longer TR, lower/variable flip angles, low-SAR refocusing.
\end{itemize}
\end{keytakeaways}

\begin{exercises}
\item How does SAR change from $1.5\ \text{T}$ to $3\ \text{T}$ at fixed flip angle?
\item Halving the flip angle changes SAR by what factor?
\item What is the IEC normal-mode whole-body SAR limit?
\end{exercises}

\furtherreading{ICNIRP~\cite{icnirp}; IEC 60601-2-33~\cite{iec60601}.}
